Fix \(F\), \(T\), \(\varnothing\ne\mathcal K\subseteq\{(g,t):2\le g\le t\le T\}\), and \(M\). For every \(Q\in\mathcal M\), let \(P=Q_O\). Suppose the off-fold learners satisfy foldwise independence, \(\max_{f,a,s}\|\widehat p_{a,s,-f}-p_{a,s}\|_{P_X,4}=O_P(a_n)\), \(\max_{f,a,s,k}\|\widehat m_{a,s,k,-f}-m_{a,s,k}\|_{P_X,4}=O_P(b_n)\), \(a_n\to0\), \(b_n\to0\), and \(a_nb_n=o(n^{-1/2})\). These learner sequences and rates are supplied as theorem-level conditional hypotheses; no uniform learner-achievability claim over \(\mathcal M\) is made. No fitted-regression envelope is separately assumed, because it follows from the primitive outcome fourth moment, overlap, and the regression error rate. Use the deterministic ridge \(\lambda_n=n^{-1/4}>0\). The score expansion and covariance lemmas then derive, rather than assume, \[\sqrt n\{\widehat\tau_{\mathrm{joint}}-\Theta_\omega^{\mathrm{agg}}(P)\}=\frac{1}{\sqrt n}\sum_{i=1}^n(b^{\mathrm{joint}})^{\prime}\Psi(O_i)+o_P(1)\rightsquigarrow N(0,V_{\mathrm{joint}}),\qquad \|\widehat\Omega-\Omega\|_{\mathrm{op}}=o_P(1).\] Moreover, \(b^{\mathrm{joint}}\) is the unique minimizer of \(b^{\prime}\Omega b\) over constant \(b\in\mathbb R^M\) satisfying \(B^{\prime}b=\omega\), \(\widehat b^{\mathrm{joint}}\to_Pb^{\mathrm{joint}}\), \(\widehat V_{\mathrm{joint}}\to_PV_{\mathrm{joint}}>0\), and \[\lim_{n\to\infty}P\{\tau_\omega(Q)\in\mathcal C_{n,1-\alpha}\}=1-\alpha.\] These are efficiency and inference claims only within the finite constant span of the retained pairwise scores, not a semiparametric efficiency bound for the full conditional-moment model. If a target coordinate has \(g=2\), its score and fitted score use \(\Delta Y_{i,k}=Y_{it}-Y_{i1}\). If \(\mathcal K=\{k\}\), the joint loading, estimator, and variance reduce to the published within-\(k\) comparison-cohort GMM objects. The zero-variance convention in Definition \(\ref{def:wald-interval}\) totalizes the finite-sample interval and is asymptotically inactive.